\section{Discussion and Conclusion}
\label{sec:discussion}
CTA investigates the representation a visual planner should predict when selecting among action chunks. A future-supervised conditional segment code improves geometry ranking on identical development banks, including over a parameter-matched direct scorer. This supports the proposed target as a useful design direction. It does not yet show that conditioning, intermediate samples, quantization, or future supervision independently causes the gain: equal-capacity conditionality/path controls and a factorized direct comparator are missing.

The evidence also exposes where the current pipeline falls short. The $L=15$ configuration shows a large source-to-prediction gap; dense geometry supervision and native success have different definitions, although the geometry oracle captures nearly all observed crossings; rare decisive events warrant focused forecast evaluation; and numerical proposal changes can alter episode outcomes. These findings motivate focused improvements in predictor/readout training, native-objective alignment, and deterministic evaluation. Their relative importance should be diagnosed within each configuration.

Practical impact depends on more than representation size. CTA uses privileged simulator labels and branched futures during training, a large frozen visual encoder, and a competent pretrained proposal policy. It cannot select an action absent from the proposal bank. Its 48-scalar deployment interface is compact, but its reader still processes context and goal tokens. Policy sampling dominates the measured workload; representation reduction is not equivalent to total compute reduction. Goal views share one PushT pose, so broad task or goal transfer remains untested.

The present conclusions are limited to reused development roots and one training seed. Historical control comparisons do not establish a reproducible learned-arm advantage, and no sealed multi-seed benchmark is claimed. The highest-priority validation is canonical closed-loop evaluation of matched continuation controls on independent roots and seeds, followed by comparable end-to-end timing and the ablations needed to attribute the ranking gain. Tasks with consequential intermediate events would separately test what the segment code preserves beyond endpoint ranking.

The central finding is therefore precise: learning a compact future-segment target can improve candidate ranking, but source quality, forecast fidelity, native-objective alignment, and reproducible control must be evaluated separately. Establishing their joint quality--compute benefit is the next step toward a stronger general planning contribution.
